% !TeX root = ../main.tex

\begin{acknowledgements}

    四年本科生活转眼即逝，从初入校园时的懵懂到今日完成这篇毕业论文，回首走过的路，所学远不止于课本中的知识与技能，更包括独立面对未知问题的方法、面对挫折时的耐心，以及科研工作所必需的求实与诚实。一段段成长的背后，离不开师长的悉心指导、朋友的相互扶持与家人的默默守候，谨借这一隅篇幅，向他们郑重致谢。

    首先要感谢我的指导老师王家祺老师。在课题选择、思路搭建、实验执行与论文写作的每一个阶段，王老师都给予了清晰而精准的指导。每当我陷入细节的纠缠，他总能以更高的视角点出问题的关键，让我得以从局部跳出、回到主线；每当我对结果的解释失之过度时，他又会及时提醒我对数据保持敬畏、对结论保持克制。他严谨求实的治学态度、对学生独立思考的尊重、以及在繁忙工作中仍坚持逐字逐句审阅论文的耐心，都是我科研生涯起点处最珍贵的榜样。在此谨向王老师致以崇高的敬意与最诚挚的感谢。

    申请季那段焦灼难捱的日子里，要感谢中山大学电子与通信工程学院的刘瀚文同学始终与我同行。我们在彼此的不确定中相互托底，把许多原本难以独自消化的压力分担成了两个人的从容。这份在最不轻松的时节里仍愿意并肩的友谊，是我本科四年最幸运的相遇之一。

    最后，要把最深的感谢留给我的家人。父母用日复一日的辛劳支撑起我无忧的求学条件，更用言传身教教会我何为正直、何为坚持；每当我面对选择犹豫不决，他们总是耐心倾听、给出最中肯的建议，又在我决定之后给予最坚定的支持。也要感谢我的哥哥，作为同样走过这条路的过来人，他在许多关键节点上给了我最踏实的提醒与陪伴，让我在前行时多了一份从容。这份家人的恩情远非言语可以承载，唯愿在未来的日子里以更踏实的脚步去回报。

    谨以此文，献给所有曾在这段旅程中给予我帮助与陪伴的人。

\end{acknowledgements}
